% 2026-10-16 (user): moved from Methods (05_methods.tex). Label sec:rh_methods_tests kept.
% The results section footnotes this analysis.

\section{Statistical tests}
\label{sec:rh_methods_tests}
\begin{itemize}
\item \textbf{Diebold--Mariano test.} To compare the predictive accuracy of two competing forecasts $A$ and $B$, we apply the \citet{DieboldMariano1995} test to the daily QLIKE loss differentials, incorporating the small-sample adjustment of \citet{HarveyLeybourneNewbold1997} for one-step-ahead forecasts:
  \[
    d_t = L_t^{A} - L_t^{B}, \qquad
    \DM = \sqrt{\frac{|\mathcal O| - 1}{|\mathcal O|}}\;
      \frac{\bar d}{\sqrt{\hat\Omega / |\mathcal O|}} .
  \]
  % 2026-10-15 (finish): parked -- the horizon symbol h (also the boosting Hessian h_{t'}) dropped; every forecast
  % is one step ahead, so h = 1 is substituted: (|O| + 1 - 2h + h(h-1)/|O|)/|O| = (|O| - 1)/|O|. Was:
  % \DM = \sqrt{\frac{|\mathcal O| + 1 - 2h + h(h-1)/|\mathcal O|}{|\mathcal O|}}\;
  %   \frac{\bar d}{\sqrt{\hat\Omega / |\mathcal O|}}, \qquad h = 1 .
  % 2026-10-15 (notation pass): parked -- the uncorrected and corrected statistics merged into one symbol, \DM
  % (the HLN-corrected one, as the code reports); n scored days -> |O|. Old display:
  % \mathrm{DM} = \frac{\bar d}{\sqrt{\hat\Omega / n}}, \qquad
  % \mathrm{DM}^{\ast} = \mathrm{DM} \sqrt{\frac{n + 1 - 2h + h(h-1)/n}{n}}, \qquad h = 1 .
  \begin{itemize}
  \item Here $A$ and $B$ denote competing forecast series $f_{m,\mathcal X,\mathcal I}$, indexed by the model class $m$, the predictor set $\mathcal X \in \{\mathrm{HAR}, \mathrm{Full}\}$, and the estimation interval set $\mathcal I \in \{\{c\}, \{1, \dots, M\}\}$.
  % 2026-10-15 (user): parked -- "out-of-sample days"; the evaluation set and the option-data subsample. Was:
  % "\item $\mathcal O$: the out-of-sample days of the exercise, numbered $t = 1, \dots, |\mathcal O|$ in time order;
  %   $\mathcal O = \mathcal T_{\mathrm{oos}}$, $|\mathcal T_{\mathrm{oos}}| = T_{\mathrm{oos}}$, in the headline
  %   comparison."
  \item The evaluation set $\mathcal O$, indexed chronologically by $t = 1, \dots, |\mathcal O|$, equals the full evaluation span $\mathcal O = \mathcal T_{\mathrm{oos}}$ ($|\mathcal T_{\mathrm{oos}}| = T_{\mathrm{oos}}$) for statistical accuracy comparisons and the option-data subsample $\mathcal T_{\mathrm{opt}} \subset \mathcal T_{\mathrm{oos}}$ for straddle return comparisons (Section~\ref{sec:rh_data}).
  \item We denote by $\bar d$ the sample mean of $d_t$ over $\mathcal O$ and compute two-sided $p$-values under the standard normal distribution as $p = 2\bigl\{1 - \Phi(|\DM|)\bigr\}$.
  \item Under this sign convention, $\DM < 0$ indicates that forecast $A$ achieves a lower average loss than forecast $B$. Throughout the paper, ``statistically significant'' refers to significance at the 5\% level ($p < 0.05$).
  % 2026-10-15 (user): merged from Forecast evaluation (moved from 06_results.tex): the p-value formula, the sign
  % convention and the meaning of "significant". Was: "\item $\bar d$: mean loss differential over $\mathcal O$;
  % two-sided $p$-value from the standard normal."
  % 2026-10-15 (finish): parked -- h dropped (one-step-ahead forecasts). Was:
  % \item $\bar d$: mean loss differential over $\mathcal O$; $h$: forecast horizon; two-sided $p$-value from the
  %   standard normal.
  \end{itemize}
% 2026-10-15 (consolidation): parked -- lag truncation q -> q_NW (q_t is the position), autocovariance lag k -> j (k
% is the HAR horizon), typed q = 6 / q = 14 -> \fqsDMLag / \fqsHACLag, plain-text citations -> \citep. The 1/|O|
% divisor of hat gamma_j is unchanged (FLAG (4) below awaits the user). Old bullets:
% 2026-10-18 (user, E13): FLAG (4) resolved -- dm_test now divides gamma_j by |O| (src/evaluation/diebold_mariano.py:67-68), matching this display.
% \item \textbf{Newey--West variance} (Newey and West 1987), Bartlett weights, lag truncation $q$:
%   \hat\Omega = \hat\gamma_0 + 2 \sum_{k=1}^{q} (1 - k/(q + 1)) \hat\gamma_k, \qquad
%   \hat\gamma_k = \frac{1}{|\mathcal O|} \sum_{t=k+1}^{|\mathcal O|} (d_t - \bar d)(d_{t-k} - \bar d).
%   DM: $q = \max(h - 1, \lfloor 4 (|\mathcal O|/100)^{2/9} \rfloor)$, which gives $q = 6$ on $\mathcal T_{\mathrm{oos}}$.
% \item \textbf{HAC $t$-statistic on the daily return differential} (Newey--West): $t = \bar d / \sqrt{\hat\Omega /
%   |\mathcal O|}$, with $d_t$ now the difference between two trading strategies' daily returns, no HLN correction;
%   $q = \lfloor 1.5\, |\mathcal O|^{1/3} \rfloor$, which gives $q = 14$ on $\mathcal T_{\mathrm{oos}}$.
% \item \textbf{Model confidence set} \citep{HansenLundeNason2011}, \textbf{superior predictive ability} (Hansen 2005)
%   and \textbf{StepM} (Romano and Wolf 2005) for the comparison of predictor sets.
\item \textbf{Newey--West long-run variance.} Following \citet{NeweyWest1987}, we estimate the long-run variance $\hat\Omega$ using Bartlett kernel weights with lag truncation parameter $q_{\mathrm{NW}}$:
  \[
    \hat\Omega = \hat\gamma_0 + 2 \sum_{j=1}^{q_{\mathrm{NW}}} \Bigl(1 - \frac{j}{q_{\mathrm{NW}} + 1}\Bigr) \hat\gamma_j,
    \qquad
    \hat\gamma_j = \frac{1}{|\mathcal O|} \sum_{t=j+1}^{|\mathcal O|} (d_t - \bar d)(d_{t-j} - \bar d).
  \]
% 2026-10-18 (user, E13): parked -- resolved; dm_test now divides gamma_j by |O|, matching the display; DM statistics regenerated on 2026-10-18 for the accuracy tables and macros (make_rehaul_results.py, make_rehaul_economics.py); the appendix study tables still need regenerating (note below). Was:
% \editc{orange}{E13}{Decision: change this display, or recompute the printed $\DM$.\\Problem: this display divides $\hat\gamma_j$ by $|\mathcal O|$, and every printed $\DM$ was computed by \texttt{dm\_test}, which divides by $|\mathcal O|-j$.\\
% The autocovariance $\hat\gamma_j$ in the display is the average of the products $(d_t-\bar d)(d_{t-j}-\bar d)$, and the divisor written in front of the sum is $|\mathcal O|$, the number of days in the comparison. There are only $|\mathcal O|-j$ such products, because the first $j$ days have no lag-$j$ partner. The function that produced the printed Diebold--Mariano statistics is \texttt{dm\_test}, and it takes those products with \texttt{np.mean}, which divides by the length of the array it is given, $|\mathcal O|-j$. The HAC $t$ on returns comes from \texttt{newey\_west\_t}, and that function divides by $|\mathcal O|$, so it matches this display. The accuracy-section $\DM$ numbers match \texttt{dm\_test}. On $\mathcal T_{\mathrm{oos}}$ the lag runs from 1 through $\rhoosDMLag$, and at each of those lags the two divisors differ.}
\editc{orange}{E13}{Decision: regenerate the remaining $\DM$ statistics.\\Problem: \texttt{dm\_test} now divides $\hat\gamma_j$ by $|\mathcal O|$, as displayed. The accuracy tables and macros were regenerated on 2026-10-18, but the $\DM$ statistics and significance stars in the appendix tuning and study tables are read from stored study outputs computed with the old divisor $|\mathcal O|-j$, so they predate the change and must be regenerated.}
  \begin{itemize}
  \item For the Diebold--Mariano test, we set the bandwidth according to the plug-in rule $q_{\mathrm{NW}} = \lfloor 4 (|\mathcal O|/100)^{2/9} \rfloor$ (its numerical value on $\mathcal T_{\mathrm{oos}}$ is reported in Section~\ref{sec:rh_exp_setup}).
  % 2026-10-15 (user): parked -- moved to Methods with Forecast evaluation, numbers removed (the lag value now lives in
  % Experimental setup, 06_results.tex). Was: "\item DM: $q_{\mathrm{NW}} = \lfloor 4 (|\mathcal O|/100)^{2/9}
  % \rfloor$, which gives $q_{\mathrm{NW}} = \fqsDMLag$ on $\mathcal T_{\mathrm{oos}}$."
  % 2026-10-15 (finish): parked -- h = 1 substituted (max(h - 1, .) = max(0, .) is the floor term). Was:
  % \item DM: $q_{\mathrm{NW}} = \max\bigl(h - 1, \lfloor 4 (|\mathcal O|/100)^{2/9} \rfloor\bigr)$, which gives
  %   $q_{\mathrm{NW}} = \fqsDMLag$ on $\mathcal T_{\mathrm{oos}}$.
  \end{itemize}
  % src/evaluation/diebold_mariano.py:33-93.
  % 2026-10-15 (notation pass): FLAG, math not changed -- dm_test computes gamma_k as np.mean(dc[k:] * dc[:-k]),
  % i.e. divides by |O| - k, not |O|; newey_west_t (the return test below) divides by |O|, as displayed.
  % 2026-10-15 (consolidation): FLAG (4) still open (user decision); the lag index is now j, the divisor unchanged.
  % 06_results.tex displays the code's |O| - j divisor for the DM statistic; the two agree once the user decides.
  % 2026-10-15 (user): the 06_results.tex display (Forecast evaluation, "Long-run variance", T_oos - j divisor and
  % explicit Bartlett weights w_j) moved to Methods and was merged into this bullet; it is parked in 06_results.tex.
  % This display (1/|O|) is now the paper's only one; FLAG (4) still decides the divisor.
  % 2026-10-18 (user, E13): FLAG (4) resolved -- dm_test now divides gamma_j by |O| (src/evaluation/diebold_mariano.py:67-68), matching this display.
% 2026-10-15 (consolidation): FLAG (7) resolved -- newey_west_t is also applied to a single strategy's mean return
% (Economic value, Table rh_econ_control), so the bullet is worded for any daily return series.
% 2026-10-21 (user, decision #2): parked -- the Newey--West t of a mean return at lag floor(1.5 |O|^(1/3)), with
% normal critical values; no published source uses the constant 1.5, so the paper no longer presents it as a rule.
% Replaced by the LLSW (2018) test with fixed-b critical values (means) and the Ledoit--Wolf (2008) test (Sharpe
% ratios), writeup/hac_inference.py. Was:
% \item \textbf{HAC $t$-statistic for mean daily returns.} For any daily return series $r_t$---either the return of an individual trading strategy or the return differential $\Delta_t$ between two strategies---we compute the unadjusted Newey--West $t$-statistic
%   \[
%     t = \frac{\bar r}{\sqrt{\hat\Omega_r / |\mathcal O|}},
%   \]
%   where $\hat\Omega_r$ is the Newey--West variance estimator above evaluated on $r_t$ in place of $d_t$, using lag truncation
  % $q_{\mathrm{NW}} = \lfloor 1.5\, |\mathcal O|^{1/3} \rfloor$ (its value on $\mathcal T_{\mathrm{oos}}$:  % 2026-10-15:
  % Experimental setup, Section~\ref{sec:rh_exp_setup}).  % parked -- the returns are on the option-data subsample
%   $q_{\mathrm{NW}} = \lfloor 1.5\, |\mathcal O|^{1/3} \rfloor$ (its numerical value on $\mathcal T_{\mathrm{opt}}$ is given in Section~\ref{sec:rh_econ_measures}).
  % 2026-10-15 (user): parked -- lag value removed from Methods, as for the DM lag (Forecast evaluation moved here,
  % numbers removed); it now lives in Experimental setup. Was: "... $q_{\mathrm{NW}} = \lfloor 1.5\, |\mathcal O|^{1/3}
  % \rfloor$, which gives $q_{\mathrm{NW}} = \fqsHACLag$ on $\mathcal T_{\mathrm{oos}}$."
  % 2026-10-15 (finish): parked -- same content; the formula moved to a display (the inline version left an underfull
  % line). Was:
  % \item \textbf{HAC $t$-statistic of a mean daily return} (Newey--West, no HLN correction):
  %   $t = \bar r / \sqrt{\hat\Omega_r / |\mathcal O|}$, with $\hat\Omega_r$ the Newey--West variance above computed on a
  %   daily return series $r_t$ in place of $d_t$: a single strategy's return (a test of its mean return) or the difference
  %   between two strategies' returns (a test of the mean return differential);
  %   $q_{\mathrm{NW}} = \lfloor 1.5\, |\mathcal O|^{1/3} \rfloor$, which gives $q_{\mathrm{NW}} = \fqsHACLag$ on
  %   $\mathcal T_{\mathrm{oos}}$.
  % notebooks/atm_straddle_lib.py:1240-1261 (newey_west_t).
\item \textbf{Test of a mean daily return.} For any daily return series $r_t$---either the return of an individual trading strategy or the return differential $\Delta_t$ between two strategies---observed on the $T = \rheSzDays$ post-warmup days $\mathcal T_{\mathrm{opt}}^{+}$ (Section~\ref{sec:rh_econ_measures}), we test whether its mean is zero with the procedure that \citet{LazarusLewisStockWatson2018} recommend for the Newey--West estimator:
  \[
    t = \frac{\bar r}{\sqrt{\hat\Omega_r / T}},
    \qquad
    \hat\Omega_r = \sum_{|j| < S} \Bigl(1 - \frac{|j|}{S}\Bigr) \hat\gamma_j,
    \qquad
    \hat\gamma_j = \frac{1}{T} \sum_{t=|j|+1}^{T} (r_t - \bar r)(r_{t-|j|} - \bar r),
  \]
  with the Bartlett kernel, bandwidth $S = \lceil 1.3\, T^{1/2} \rceil = \rheHACS$, and fixed-$b$ critical values for $b = S/T$ \citep{KieferVogelsang2005}.
  \begin{itemize}
  \item The autocovariances divide by $T$, with no small-sample factor; at lag truncation $S - 1$ the estimator equals the Newey--West variance of the previous item.
  \item This bandwidth is larger than that of conventional rules such as $0.75\, T^{1/3}$, which \citet{LazarusLewisStockWatson2018} obtain from the formula of \citet{Andrews1991} for a first-order autoregression with coefficient $0.25$. Combined with fixed-$b$ critical values, it reduces the size distortion of the test when the returns are serially correlated \citep{LazarusLewisStockWatson2018}.
  \item Under the null hypothesis and fixed $b$, the limiting distribution of $t$ is a functional of Brownian motion that depends only on the kernel and on $b$ \citep{KieferVogelsang2005}. We obtain its quantiles by simulation, computing the same statistic on 1{,}000{,}000 independent series of 1{,}000 i.i.d.\ standard normal draws with bandwidth $1{,}000\,b$. This gives the two-sided 5\% critical value \rheHACCV{} and the $p$-values reported in the text, in place of the standard normal value 1.96.
  \end{itemize}
\item \textbf{Test of equal Sharpe ratios.} A test of $\bar\Delta = 0$ concerns the mean return differential, not the difference between two Sharpe ratios, which divides each strategy's mean return by its own standard deviation. For every statement about a Sharpe ratio difference we therefore use the studentized circular block bootstrap test of \citet{LedoitWolf2008}.
  \begin{itemize}
  \item With $u = \bigl(\mu_A, \mu_B, \mathrm E[(r^A_t)^2], \mathrm E[(r^B_t)^2]\bigr)$, the difference of the Sharpe ratios of the daily returns, $\mu_A/\sigma_A - \mu_B/\sigma_B$, is a smooth function $f(u)$; its standard error follows from the delta method and the long-run covariance matrix $\Psi$ of $\bigl(r^A_t, r^B_t, (r^A_t)^2, (r^B_t)^2\bigr)$.
  \item The test statistic is $d = |f(\hat u)| / \mathrm{se}(\hat u)$. Each of \rheLWBoot{} bootstrap samples joins blocks of \rheLWBlock{} consecutive days, drawn with replacement from the $T$ post-warmup days arranged on a circle, and gives $d^{*} = |f(\hat u^{*}) - f(\hat u)| / \mathrm{se}(\hat u^{*})$, with $\Psi$ estimated in that sample by the block-based estimator of \citet{LedoitWolf2008}. The $p$-value is $\bigl(\#\{d^{*} \ge d\} + 1\bigr)/(\rheLWBoot + 1)$.
  \item In the original sample, $\Psi$ is estimated with the Bartlett kernel at bandwidth equal to the block length, whose expectation is approximately that of the block-based estimator; \citet{LedoitWolf2008} use a prewhitened quadratic-spectral kernel at this step. The block length is fixed at \rheLWBlock{} days rather than calibrated; the next item reports the number of rejections at block lengths \rheLWBlockLo{} and \rheLWBlockHi{}. The random seed is fixed.
  \item The test concerns the Sharpe ratios of daily returns over the $T$ post-warmup days, without the year-by-year annualization of Section~\ref{sec:rh_econ_measures}.
  \end{itemize}
\item \textbf{Robustness to the long-run variance estimator.} Table~\ref{tab:rh_app_hac_robust} lists every comparison of strategy returns whose statistical significance Section~\ref{sec:rh_econ} or Appendix~\ref{sec:rh_econ_last30} states. For each comparison it reports the $t$-statistic of the mean return differential with the Newey--West variance at lag truncation $q_{\mathrm{NW}} = \lfloor 4 (T/100)^{2/9} \rfloor = \rhrRobLagNW$ \citep{NeweyWest1994} and at $q_{\mathrm{NW}} = \lfloor 1.5\, T^{1/3} \rfloor = \rhrRobLagOld$, the lag truncation of an earlier version of this paper, both with standard normal $p$-values; the fixed-$b$ test above; and the $p$-value of the \citet{LedoitWolf2008} test. Statements in the text about mean returns use the fixed-$b$ test, and statements about Sharpe ratios use the \citet{LedoitWolf2008} test.
  \begin{itemize}
  \item Of the \rhrRobN{} tests of a mean, \rhrRobFbSig{} are statistically significant at the 5\% level under the fixed-$b$ test, \rhrRobNWFiveSig{} at lag truncation \rhrRobLagNW{} and \rhrRobNWTwelveSig{} at lag truncation \rhrRobLagOld{} with standard normal critical values.
  \item Of the \rhrRobNPairs{} tests of equal Sharpe ratios, \rhrRobLWSig{} reject at the 5\% level with block length \rheLWBlock{}, \rhrRobLWSigLo{} with block length \rheLWBlockLo{} and \rhrRobLWSigHi{} with block length \rheLWBlockHi{}.
  \end{itemize}
% generated by writeup/make_rehaul_results.py (2026-10-21, decision #2): the tests of strategy returns behind
% the significance statements of Section 7 and Appendix F, on the 614 post-warmup days, under four procedures.
% 2026-10-26 (D19): Sharpe-ratio differences by Ledoit-Wolf (the paper's test); the bootstrap column is parked.
% Do not edit by hand.
\begin{table}[tbp]
\centering
\scriptsize
\setlength{\tabcolsep}{3.5pt}
\begin{threeparttable}
\caption{Tests of strategy returns under alternative long-run variance estimators}
\label{tab:rh_app_hac_robust}
\begin{tabular}{@{}lrrrrrrrrr@{}}
\toprule
 & & & \multicolumn{2}{c}{NW, $q_{\mathrm{NW}} = \rhrRobLagNW$} & \multicolumn{2}{c}{NW, $q_{\mathrm{NW}} = \rhrRobLagOld$} & \multicolumn{2}{c}{NW, $S = \rheHACS$, fixed-$b$} & LW \\
\cmidrule(lr){4-5}\cmidrule(lr){6-7}\cmidrule(lr){8-9}\cmidrule(l){10-10}
 & $\bar\Delta$ (\%) & $\Delta\SR$ & $t$ & $p$ & $t$ & $p$ & $t$ & $p$ & $p$ \\
\midrule
\multicolumn{10}{@{}l}{\textit{Panel A: Pooled intraday model minus short-straddle benchmark}} \\
\addlinespace[2pt]
Ridge, HAR & $4.76$ & $0.63$ & $0.93$ & $0.350$ & $0.95$ & $0.344$ & $0.90$ & $0.388$ & $0.350$ \\
Ridge, full & $9.89$ & $1.32$ & $1.88$ & $0.060$ & $1.82$ & $0.069$ & $1.70$ & $0.110$ & $0.059$ \\
Lasso, HAR & $4.76$ & $0.63$ & $0.93$ & $0.350$ & $0.95$ & $0.344$ & $0.90$ & $0.388$ & $0.350$ \\
Lasso, full & $10.52$ & $1.40$ & $2.03$ & $0.042$ & $2.01$ & $0.045$ & $1.93$ & $0.071$ & $0.039$ \\
Elastic net, HAR & $4.76$ & $0.63$ & $0.93$ & $0.350$ & $0.95$ & $0.344$ & $0.90$ & $0.388$ & $0.350$ \\
Elastic net, full & $11.46$ & $1.53$ & $2.22$ & $0.026$ & $2.13$ & $0.033$ & $1.95$ & $0.068$ & $0.026$ \\
LightGBM, HAR & $2.50$ & $0.33$ & $0.48$ & $0.634$ & $0.52$ & $0.606$ & $0.53$ & $0.610$ & $0.614$ \\
LightGBM, full & $9.11$ & $1.21$ & $1.76$ & $0.078$ & $1.76$ & $0.078$ & $1.68$ & $0.113$ & $0.070$ \\
XGBoost, HAR & $4.96$ & $0.66$ & $0.92$ & $0.356$ & $1.02$ & $0.308$ & $1.08$ & $0.304$ & $0.315$ \\
XGBoost, full & $9.17$ & $1.22$ & $1.77$ & $0.076$ & $1.76$ & $0.078$ & $1.61$ & $0.127$ & $0.064$ \\
Random forest, HAR & $-0.03$ & $0.00$ & $-0.01$ & $0.994$ & $-0.01$ & $0.993$ & $-0.01$ & $0.993$ & $0.992$ \\
Random forest, full & $11.64$ & $1.56$ & $2.28$ & $0.023$ & $2.33$ & $0.020$ & $2.14$ & $0.046$ & $0.018$ \\
\addlinespace[4pt]
\multicolumn{10}{@{}l}{\textit{Panel B: Pooled intraday model, full predictor set minus HAR benchmark}} \\
\addlinespace[2pt]
Ridge & $5.13$ & $0.69$ & $1.43$ & $0.153$ & $1.37$ & $0.171$ & $1.32$ & $0.209$ & $0.171$ \\
Lasso & $5.77$ & $0.77$ & $1.60$ & $0.110$ & $1.54$ & $0.122$ & $1.59$ & $0.132$ & $0.120$ \\
Elastic net & $6.70$ & $0.90$ & $1.91$ & $0.056$ & $1.79$ & $0.074$ & $1.67$ & $0.115$ & $0.074$ \\
LightGBM & $6.61$ & $0.88$ & $1.99$ & $0.047$ & $1.94$ & $0.053$ & $1.89$ & $0.076$ & $0.050$ \\
XGBoost & $4.21$ & $0.56$ & $1.27$ & $0.203$ & $1.27$ & $0.204$ & $1.16$ & $0.268$ & $0.204$ \\
Random forest & $11.68$ & $1.56$ & $2.51$ & $0.012$ & $2.56$ & $0.010$ & $2.42$ & $0.026$ & $0.020$ \\
\addlinespace[4pt]
\multicolumn{10}{@{}l}{\textit{Panel C: Pooled intraday minus closing-interval model}} \\
\addlinespace[2pt]
Ridge, HAR & $-4.43$ & $-0.59$ & $-1.52$ & $0.128$ & $-1.55$ & $0.121$ & $-1.36$ & $0.197$ & $0.134$ \\
Ridge, full & $1.62$ & $0.22$ & $0.54$ & $0.592$ & $0.57$ & $0.569$ & $0.58$ & $0.579$ & $0.584$ \\
Lasso, HAR & $-4.43$ & $-0.59$ & $-1.52$ & $0.128$ & $-1.55$ & $0.121$ & $-1.36$ & $0.197$ & $0.134$ \\
Lasso, full & $0.44$ & $0.06$ & $0.12$ & $0.901$ & $0.13$ & $0.898$ & $0.13$ & $0.900$ & $0.891$ \\
Elastic net, HAR & $-4.43$ & $-0.59$ & $-1.52$ & $0.128$ & $-1.55$ & $0.121$ & $-1.36$ & $0.197$ & $0.134$ \\
Elastic net, full & $0.22$ & $0.03$ & $0.07$ & $0.944$ & $0.07$ & $0.943$ & $0.07$ & $0.948$ & $0.947$ \\
LightGBM, HAR & $1.10$ & $0.15$ & $0.26$ & $0.792$ & $0.27$ & $0.790$ & $0.26$ & $0.804$ & $0.813$ \\
LightGBM, full & $2.17$ & $0.29$ & $0.51$ & $0.607$ & $0.52$ & $0.602$ & $0.56$ & $0.592$ & $0.639$ \\
XGBoost, HAR & $1.59$ & $0.21$ & $0.64$ & $0.522$ & $0.67$ & $0.502$ & $0.67$ & $0.519$ & $0.519$ \\
XGBoost, full & $3.32$ & $0.44$ & $0.79$ & $0.432$ & $0.82$ & $0.412$ & $0.79$ & $0.447$ & $0.445$ \\
Random forest, HAR & $-1.01$ & $-0.13$ & $-0.40$ & $0.690$ & $-0.41$ & $0.680$ & $-0.46$ & $0.659$ & $0.678$ \\
Random forest, full & $8.08$ & $1.08$ & $2.68$ & $0.007$ & $2.91$ & $0.004$ & $2.70$ & $0.013$ & $0.007$ \\
\addlinespace[4pt]
\multicolumn{10}{@{}l}{\textit{Panel D: Closing-interval model, full predictor set minus HAR benchmark}} \\
\addlinespace[2pt]
Elastic net & $2.05$ & $0.28$ & $0.78$ & $0.436$ & $0.78$ & $0.437$ & $0.74$ & $0.476$ & $0.450$ \\
Lasso & $0.89$ & $0.12$ & $0.35$ & $0.727$ & $0.34$ & $0.733$ & $0.35$ & $0.739$ & $0.737$ \\
\addlinespace[4pt]
\multicolumn{10}{@{}l}{\textit{Panel E: Closing-interval model minus short-straddle benchmark}} \\
\addlinespace[2pt]
Elastic net, full & $11.24$ & $1.50$ & $1.95$ & $0.051$ & $2.00$ & $0.046$ & $1.92$ & $0.072$ & $0.035$ \\
\addlinespace[4pt]
\multicolumn{10}{@{}l}{\textit{Panel F: Short-straddle benchmark, mean return}} \\
\addlinespace[2pt]
Short straddle & $2.43$ & -- & $0.58$ & $0.559$ & $0.62$ & $0.536$ & $0.67$ & $0.522$ & -- \\
\bottomrule
\end{tabular}
\begin{tablenotes}[flushleft]
\footnotesize
\item \textit{Notes:} Each row is a comparison whose statistical significance the text states, on the \rheSzDays{} post-warmup days of $\mathcal T_{\mathrm{opt}}$. $\bar\Delta$ is the mean daily return differential at mid-quotes in percent of the straddle premium (the mean return itself in the last panel), and $\Delta\SR$ the difference of the annualized Sharpe ratios of Section~\ref{sec:rh_econ_measures}. The $t$-statistics test $\bar\Delta = 0$ with the Newey--West long-run variance: at lag truncation $q_{\mathrm{NW}} = \lfloor 4(T/100)^{2/9}\rfloor = \rhrRobLagNW$ \citep{NeweyWest1994} and at $q_{\mathrm{NW}} = \lfloor 1.5\,T^{1/3}\rfloor = \rhrRobLagOld$, both with standard normal $p$-values, and at bandwidth $S = \lceil 1.3\,T^{1/2}\rceil = \rheHACS$ with fixed-$b$ $p$-values \citep{LazarusLewisStockWatson2018,KieferVogelsang2005}, $T = \rheSzDays$. LW is the $p$-value of the \citet{LedoitWolf2008} test of equal Sharpe ratios $\SR$ (circular block bootstrap, block length \rheLWBlock{}, \rheLWBoot{} resamples). The text uses the fixed-$b$ test for statements about mean returns and the LW test for statements about Sharpe ratios. All $p$-values are two-sided.
\end{tablenotes}
\end{threeparttable}
\end{table}

\item \textbf{Multiple-comparison procedures.} For joint comparisons across predictor sets and model families, one may also consider the model confidence set of \citet{HansenLundeNason2011}, the superior predictive ability test of \citet{Hansen2005}, and the StepM procedure of \citet{RomanoWolf2005}.
\editc{blue}{E14}{Decision: drop the sentence, or report the set.\\Problem: the model confidence set, the superior predictive ability test, and StepM are named, and the paper reports none of their outputs.\\
This bullet names three procedures and says they are for the comparison of predictor sets: the model confidence set, which returns the set of models not beaten at a stated confidence level; the superior predictive ability test, which returns a $p$-value against a benchmark; and StepM, which steps through the models and rejects some of them. No table in the paper states an MCS membership, an SPA $p$-value, or a StepM rejection. The next bullet says the main tables are point estimates and contain no bootstrap intervals, which is the kind of output those three procedures produce. The comparisons that are printed are Diebold--Mariano statistics on QLIKE and on squared error.}
% 2026-10-21 (user, decision #2): parked -- "Diebold--Mariano or HAC $t$-statistics". Was: "\item \textbf{Reporting convention.} All main tables report point estimates accompanied by pairwise Diebold--Mariano or HAC $t$-statistics, without bootstrap confidence intervals."
\item \textbf{Reporting convention.} All main tables report point estimates accompanied by pairwise Diebold--Mariano statistics, $t$-statistics of mean returns with fixed-$b$ critical values, or \citet{LedoitWolf2008} bootstrap $p$-values for Sharpe ratio differences, without bootstrap confidence intervals.
\end{itemize}

